% Generated draft: requires mathematical and source review.
\documentclass[11pt]{article}
\usepackage[margin=25mm]{geometry}
\usepackage{fontspec}
\setmainfont{Latin Modern Roman}
\usepackage{amsmath,amssymb,mathtools,unicode-math}
\setmathfont{Latin Modern Math}
\usepackage{enumitem,needspace,xurl,hyperref}
\setlength{\parindent}{0pt}
\setlength{\parskip}{5pt}
\setlength{\emergencystretch}{3em}
\newcommand{\N}{\mathbb{N}}
\newcommand{\Z}{\mathbb{Z}}
\newcommand{\Q}{\mathbb{Q}}
\newcommand{\R}{\mathbb{R}}
\newcommand{\C}{\mathbb{C}}
\newcommand{\F}{\mathbb{F}}
\newcommand{\A}{\mathbb{A}}
\newcommand{\PP}{\mathbb{P}}
\newcommand{\E}{\mathbb{E}}
\newcommand{\eps}{\varepsilon}
\newcommand{\dd}{\,\mathrm{d}}
\begin{document}
\section*{UnsolvedMath ID: 30005380}
\textbf{Source problem code:} catalogue entry
\subsection{Definitions and mathematical statement}
A real-valued quasimorphism on a group \(G\) is a function \(q:G\to\R\) for which there exists \(D\geq0\) such that
\[
|q(xy)-q(x)-q(y)|\leq D\qquad(x,y\in G).
\]
It is homogeneous if \(q(g^n)=nq(g)\) for every \(g\in G\) and \(n\in\Z\). For an additive group \(C\) equipped with a metric \(d\), a real-valued homomorphism \(I:C\to\R\) is \(1\)-Lipschitz if
\[
I(x+y)=I(x)+I(y),\qquad |I(x)-I(y)|\leq d(x,y)\qquad(x,y\in C).
\]
Here \(C\) and \(d\) are placeholders: the supplied record does not specify the concordance group, its category, or its metric.

The supplied question asks whether, for the braid quasimorphism \(f\) under discussion, there exist a \(1\)-Lipschitz concordance homomorphism \(I\) and a scalar \(r\in\R\) such that \(f=rI\). This equality requires an unspecified comparison between a concordance homomorphism and a function on braids. The background instead describes obtaining braid quasimorphisms by homogenizing knot invariants. If \(B\) denotes the intended braid group and \(\mathcal H(I):B\to\R\) denotes the intended construction from \(I\), that reading would give the schematic question
\[
\exists I\;\exists r\in\R\;\forall\beta\in B,
\qquad f(\beta)=r\,\mathcal H(I)(\beta),
\]
where \(I\) must satisfy the homomorphism and Lipschitz conditions above. Neither \(B\), the admissible class of \(f\), nor \(\mathcal H\) is defined in the supplied excerpt, so this is a conditional formulation, not a fully specified mathematical statement.

There is also an unresolved quantifier discrepancy. The question text refers to a particular \(f\), whereas the supplied literature assessment describes a single \(I\) whose homogenization generates every such braid quasimorphism. If \(\mathcal F\) were the intended class, the latter would have the different schematic form
\[
\exists I\;\forall f\in\mathcal F\;\exists r_f\in\R\;\forall\beta\in B,
\qquad f(\beta)=r_f\,\mathcal H(I)(\beta).
\]
The record does not determine which quantifier order is intended.

The adjacent source passage says that an unnamed theorem answers \([\mathrm{HKK}^{+}20,\ \text{Question }1.6]\), and that this theorem and the slice-Bennequin inequality can be proved by realizing the involved homogeneous quasimorphisms, denoted \(\mathrm{wr}\) and \(\omega\), respectively, as homogenizations of a certain type of knot invariant. The theorem, the invariants, and the required constructions are not supplied; these contextual claims do not specify or resolve the question above.
\subsection{Short English statement}
The question asks whether a braid quasimorphism arises, up to a real scalar, from a \(1\)-Lipschitz concordance homomorphism through an intended homogenization construction. The supplied record leaves the domains, metric, construction, admissible quasimorphisms, and quantifier order unresolved. It labels the problem open, but that label and the claimed source verification are supplied curation requiring human review.
\subsection{Sources}
\begin{itemize}
\item Local UnsolvedMath registry, permanent ID \(30005380\), source problem code \texttt{catalogue entry}, title: Lipschitz Concordance Homomorphisms for Braid Quasimorphisms. The record assigns proposed year 2023 and names no proposer.
\item Supplied primary source citation: \emph{Morphisms in Low Dimensions} (2023), Oberwolfach Reports, pp. 215--259, DOI \texttt{10.4171/owr/2023/4}. \url{https://doi.org/10.4171/owr/2023/4}. The background describes these reports as documenting mathematical workshops and their open-problem discussions.
\item Supplied literature source, labeled Oberwolfach Report 4/2023: \url{https://doi.org/10.4171/owr/2023/4}.
\item Supplied literature source, attributed to Feller and labeled \emph{Slice-Bennequin inequality for FDTC}: \url{https://doi.org/10.1017/S1474748024000276}.
\item The excerpt mentions \([\mathrm{HKK}^{+}20,\ \text{Question }1.6]\), but supplies no full bibliographic reference.
\item Literature assessment dated exactly \texttt{2026-08-22}: the supplied curation labels the problem open and reports that no general resolution was found for the existence of a single \(1\)-Lipschitz concordance homomorphism whose homogenization generates every such braid quasimorphism. This is supplied dated curation derived from the supplied Oberwolfach material, not an independently verified assessment. Time-sensitive status and cited claims require independent verification.
\item The registry labels the statement \texttt{corrected\_verified} and reports that it was checked against the official Oberwolfach report. This verification claim is supplied metadata; no source visit or independent mathematical verification is asserted here.
\end{itemize}
\end{document}
